% Môn: Toán
% Lớp: 10
% Chương: Chương I. Mệnh đề và tập hợp
% Bài: Bài 2. Tập hợp và các phép toán trên tập hợp · Xác định tập hợp, phần tử và tập con
% Số câu: 56
% App đọc file này trực tiếp — sửa rồi Commit trên GitHub, bấm Đồng bộ GitHub .tex

% Bài 2. Tập hợp và các phép toán trên tập hợp



% ===== Câu 4 =====
% ID: T10C1B2-01-TN
% Mức: NB
\dangbt{Xác định tập hợp, phần tử và tập con}
\begin{ex}
% ID: T10C1B2-01-TN
% Mức: NB
Cho $x$ là một phần tử của tập hợp $X$. Mệnh đề nào sau đây là $\textbf{đúng}$?
\choice
{\True $x \in X$}
{$X \in x$}
{$x \subset X$}
{$\{x\} \in X$}
\loigiai{
Mệnh đề là đúng là $x \in X$. Kí hiệu $x$ là một phần tử của tập hợp $X$.
}
\end{ex}

% ===== Câu 16 =====
% ID: T10C1B2-04-TN
% Mức: NB
\begin{ex}
% ID: T10C1B2-04-TN
% Mức: NB
Cho tập hợp $A=\{2;4;7;8;9;12\}$. Đâu là một tập hợp con của $A$?
\choice
{$\mathscr{D}=\{5\}$}
{$C=\{2;6;9\}$}
{$B=\{3;4;7\}$}
{\True $X=\{2;12\}$}
\loigiai{
tập hợp con của tập hợp $A$ là $X=\{2;12\}$.
}
\end{ex}

% ===== Câu 24 =====
% ID: T10C1B2-07-TN
% Mức: NB
\begin{ex}
% ID: T10C1B2-07-TN
% Mức: NB
Hãy liệt kê các phần tử của tập hợp $X=\left\{n \in \mathbb{N}^{*} \mid n \leq 3\right\}$.
\choice
{$X=\{1 ; 2\}$}
{\True $X=\{1 ; 2 ; 3\}$}
{$X=\{0 ; 1 ; 2\}$}
{$X=\{0 ; 1 ; 2 ; 3\}$}
\loigiai{
$X=\{1 ; 2 ; 3\}$.
}
\end{ex}

% ===== Câu 28 =====
% ID: T10C1B2-10-TN
% Mức: NB
\begin{ex}
% ID: T10C1B2-10-TN
% Mức: NB
Cho tập hợp $A=\{x \mid x \in \mathbb{N}, x \leq 6\}$. Số phần tử của tập hợp $A$ là
\choice
{$8$}
{\True $7$}
{$5$}
{$6$}
\loigiai{
$A=\{x \mid x \in \mathbb{N}, x \leq 6\} \Leftrightarrow A=\{0; 1 ; 2 ; 3 ; 4 ; 5;6 \}$. Do đó số phần tử của tập hợp $A$ là $7$.
}
\end{ex}

% ===== Câu 29 =====
% ID: T10C1B2-11-TN
% Mức: NB
\begin{ex}
% ID: T10C1B2-11-TN
% Mức: NB
Số phần tử của tập $C=\{ x \in \mathbb{Z} \mid -1 \leq x \leq 1\}$ là
\choice
{\True $3$}
{$2$}
{$1$}
{$6$}
\loigiai{
Ta có $C=\{-1 ; 0 ; 1\}$ nên $C$ có $3$ phần tử.
}
\end{ex}

% ===== Câu 30 =====
% ID: T10C1B2-12-TN
% Mức: NB
\begin{ex}
% ID: T10C1B2-12-TN
% Mức: NB
Kí hiệu nào sau đây dùng để viết đúng mệnh đề «~ $\sqrt{2}$ không phải là số hữu tỉ~»?
\choice
{$\sqrt{2}\subset \mathbb{Q}$}
{\True $\sqrt{2}\notin \mathbb{Q}$}
{$\sqrt{2}\ne \mathbb{Q}$}
{$\sqrt{2}\in \mathbb{Q}$}
\loigiai{
Cách viết đúng là $\sqrt{2} \notin \mathbb{Q}$.
}
\end{ex}

% ===== Câu 31 =====
% ID: T10C1B2-13-TN
% Mức: NB
\begin{ex}
% ID: T10C1B2-13-TN
% Mức: NB
Cách viết nào sau đây là đúng để mô tả «$x$ là một phần tử của tập hợp $A$»
\choice
{\True $x\in A$}
{$\{x\} \in A$}
{$x\subset A$}
{$\{x\} \not \subset A$}
\loigiai{
Vì $x$ là phần tử cho nên ta có $x\in A$ là cách viết đúng.
}
\end{ex}

% ===== Câu 34 =====
% ID: T10C1B2-16-TN
% Mức: NB
\begin{ex}
% ID: T10C1B2-16-TN
% Mức: NB
Cho tập hợp $X=\{0 ; 2 ; 5\}$. Tập hợp nào dưới đây $\textbf{không}$ phải là tập con của tập hợp $X$?
\choice
{$\varnothing$}
{$\{2\}$}
{$\{0 ; 2 ; 5\}$}
{\True $\{0 ; 1 ; 2\}$}
\loigiai{
Ta có $1 \notin X$, vậy $\{0 ; 1 ; 2\} \not \subset\{0 ; 2 ; 5\}$.
}
\end{ex}

% ===== Câu 39 =====
% ID: T10C1B2-21-TN
% Mức: NB
\begin{ex}
% ID: T10C1B2-21-TN
% Mức: NB
Cho tập hợp $A=\{ x \in \mathbb{N} |~(x-1)(2x+4)=0 \}$. Tập $A$ được viết theo kiểu liệt kê là
\choice
{$\{ -1;4 \}$}
{\True $\{ 1 \}$}
{$\{ 1;-2 \}$}
{$\{ -1;-2 \}$}
\loigiai{
Ta có
  (x-1)(2x+4)=0
 $\Leftrightarrow$ hoặc{&x-1=0
&2x+4=0}
 $\Leftrightarrow$ hoặc{&x=1
&x=-2 $\notin\mathbb{N}$.} 
 Vậy $A=\{ x \in \mathbb{N} |~ (x-1)(2x+4)=0 \} = \{ 1 \}$.
}
\end{ex}

% ===== Câu 47 =====
% ID: T10C1B2-29-TN
% Mức: VD
\begin{ex}
% ID: T10C1B2-29-TN
% Mức: VD
Cho tập $M=\{x \in \mathbb{R} \mid (x+3)(x-1) \geqslant 0\}$. Khi đó
\choice
{$M=(-\infty ;-1] \cup[3 ;+\infty)$}
{\True $M=(-\infty ;-3] \cup[1 ;+\infty)$}
{$M=[1 ;+\infty)$}
{$M=(-\infty ;-3]$}
\loigiai{
Ta có $(x+3)(x-1) \geqslant 0 \Leftrightarrow x \leqslant -3$ hoặc $x \geqslant 1$.
Vậy $M = (-\infty; -3] \cup [1; +\infty)$.
}
\end{ex}

% ===== Câu 48 =====
% ID: T10C1B2-30-TN
% Mức: VD
\begin{ex}
% ID: T10C1B2-30-TN
% Mức: VD
Trong các tập hợp sau, tập nào là tập rỗng?
\choice
{$A= \{ x \in \mathbb{R} |~ x^2+5x-6=0 \}$}
{$B=\{ x \in \mathbb{Q} |~ 3x^2-5x+2=0 \}$}
{\True $C=\{ x\in \mathbb{Z} |~ x^2+x-1=0 \}$}
{$D=\{ x \in \mathbb{R} |~ x^2+5x-1=0 \}$}
\loigiai{
Ta có
$x^2+5x-6=0
 \Leftrightarrow hoặc{&x=1 \in \mathbb{R}\\&x=-6 \in \mathbb{R}}$. Vậy$A=\{ -6;1 \} \neq \varnothing$.
$3x^2-5x+2=0
 \Leftrightarrow hoặc{&x=1 \in \mathbb{Q}\\&x=\dfrac{2}{3} \in \mathbb{Q}}$. Vậy$B=\left\{ 1;\dfrac{2}{3} \right\} \neq \varnothing$.
$x^2+x-1=0
 \Leftrightarrow hoặc{&x=\dfrac{-1+\sqrt{5}}{2} \notin \mathbb{Z}\\&x=\dfrac{-1-\sqrt{5}}{2} \notin \mathbb{Z}}$. Vì$x \in \mathbb{Z}$nên$C=\varnothing$.
$x^2+5x-1=0
 \Leftrightarrow hoặc{&x=\dfrac{-5+\sqrt{29}}{2} \in \mathbb{R}\\&x=\dfrac{-5-\sqrt{29}}{2} \in \mathbb{R}}$. Vậy$D=\left\{ \dfrac{-5+\sqrt{29}}{2};\dfrac{-5-\sqrt{29}}{2} \right\} \neq \varnothing$.
}
\end{ex}

% ===== Câu 51 =====
% ID: T10C1B2-33-TN
% Mức: VD
\begin{ex}
% ID: T10C1B2-33-TN
% Mức: VD
Cho tập $A=\{ x\in \mathbb{N}|~(2-x)(x^2-3x-4)=0 \}$. Hỏi tập $A$ có bao nhiêu tập con ?
\choice
{$2$}
{\True $4$}
{$7$}
{$8$}
\loigiai{
Ta có $(2-x)(x^2-3x-4)=0 \Leftrightarrow hoặc{
 & x=2 \in \mathbb{N} \\
 & x=-1 \notin \mathbb{N} \\
 & x=4 \in \mathbb{N}.}$
Vậy $A=\{ 2;4 \}$ nên $A$ có các tập con là $\{ 2\}$; $\{ 4\}$; $\{ 2;4\}$; $\varnothing$.
}
\end{ex}
